\chapter{Introduzione all'ottimizzazione}
Ogni volta che si hanno a disposizione risorse limitate e più modi di utilizzarle, ci si trova davanti a un problema di scelta: tra tutte le alternative possibili, qual è la migliore? L'ottimizzazione matematica nasce per rispondere a questa domanda in modo rigoroso, ovvero traducendo il problema reale in un modello matematico che possa essere analizzato e risolto con metodi numerici. In questo capitolo introduciamo il linguaggio di base che useremo per tutto il corso: cosa si intende per problema di ottimizzazione, come è fatto e in quali casi ammette (o non ammette) una soluzione.

\section{Un primo esempio: i mattoncini lego}\label{sec:lego-intro}
Per fissare le idee, partiamo da un problema concreto. Supponiamo di avere a disposizione 8 mattoncini piccoli e 6 mattoncini grandi, e di poter costruire due tipi di oggetti:
\begin{itemize}
    \item \textbf{Tavolo}: richiede 2 mattoncini piccoli e 2 grandi, e permette di guadagnare 20 euro.
    \item \textbf{Sedia}: richiede 2 mattoncini piccoli e 1 grande, e permette di guadagnare 15 euro.
\end{itemize}

Il problema è decidere quanti tavoli e quante sedie costruire in modo da massimizzare il profitto. Prima di poterlo risolvere servono però gli strumenti per formularlo: la formulazione matematica verrà costruita nel capitolo~\ref{ch:pl}.

\section{Cos'è l'ottimizzazione}\label{sec:cos-e}
L'\textbf{ottimizzazione matematica}, detta anche \textbf{programmazione matematica}, è una branca della \textbf{Ricerca Operativa}, ovvero la disciplina che si occupa di formulare e risolvere modelli matematici per processi decisionali complessi (per questo motivo la Ricerca Operativa è detta anche la \emph{scienza delle decisioni}). In particolare, l'ottimizzazione si occupa di problemi decisionali che sono:
\begin{itemize}
    \item \textbf{Con un solo decisore}: la scelta è compiuta da un unico soggetto.
    \item \textbf{Con un solo obiettivo}: esiste un unico criterio che permette di confrontare le diverse alternative.
    \item \textbf{Deterministici}: i dati si considerano certi, senza alcuna componente probabilistica.
\end{itemize}

È importante precisare che i problemi con più decisori, ognuno con i propri obiettivi, non rientrano in questo ambito: la loro risoluzione è studiata dalla \textbf{teoria dei giochi}.

% TODO: problema di Dantzig, da leggere nelle slide del corso.

\section{Approccio modellistico}
Il passaggio da un problema reale alla sua soluzione segue, in generale, cinque fasi:
\begin{enumerate}
    \item \textbf{Formulazione del problema}: si individuano gli obiettivi, si determinano le restrizioni e si raccolgono i dati.
    \item \textbf{Costruzione del modello}: si cerca un compromesso tra precisione e trattabilità matematica, individuando le variabili, la funzione obiettivo e i vincoli.
    \item \textbf{Analisi del modello}: si individuano i metodi numerici adatti e si costruisce l'algoritmo risolutivo.
    \item \textbf{Soluzione numerica}: si esegue l'algoritmo e si calcola la soluzione.
    \item \textbf{Validazione del modello}: si verifica la soluzione ottenuta (collaudo) e, se necessario, si migliora il modello.
\end{enumerate}

Un modello troppo dettagliato rischia di essere impossibile da risolvere, mentre un modello troppo semplice potrebbe non descrivere abbastanza bene il problema reale: per questo motivo l'ultima fase può riportare indietro alle prime, e il processo viene spesso ripetuto più volte.

\section{Com'è fatto un problema di ottimizzazione}
La forma più generale di un problema di ottimizzazione è:
\begin{equation}\label{eq:opt-generale}
\min f(x), \quad x \in X, \quad f: \mathbb{R}^n \to \mathbb{R}
\end{equation}
dove:
\begin{itemize}
    \item $f(x)$ è la \textbf{funzione obiettivo}, ovvero il criterio di scelta tra le diverse alternative.
    \item $X \subseteq \mathbb{R}^n$ è la \textbf{regione ammissibile}, ovvero l'insieme dei punti che rispettano i vincoli del problema; un punto $x \in X$ è detto \textbf{punto ammissibile}.
    \item $x = (x_1, \dots, x_n)^T$ è il vettore delle \textbf{variabili}.
\end{itemize}

Un problema di ottimizzazione consiste quindi nel determinare, se esiste, un punto di minimo della funzione $f$ tra i punti dell'insieme $X$.

\begin{osservazione}[Problemi di massimo e di minimo]
Nel seguito parleremo principalmente di problemi di minimo, perché un problema di massimo può sempre essere ricondotto a uno di minimo cambiando di segno la funzione obiettivo. Infatti i punti di massimo di $f$ su $X$ coincidono con i punti di minimo di $-f$, e di conseguenza:
\[
\max_{x \in X} f(x) = -\min_{x \in X} \big(-f(x)\big).
\]
\end{osservazione}

\paragraph{Esempio.}
Consideriamo la funzione
\[
f(x) = -x^2 + 10x
\]
ovvero una parabola con vertice $V_0 = (5, 25)$, rappresentata nella figura~\ref{fig:parabole} a sinistra. Il massimo di $f(x)$ è l'ordinata del vertice, cioè $25$. Se consideriamo invece $-f(x) = x^2 - 10x$, la parabola si ribalta (figura~\ref{fig:parabole}, a destra) e il suo minimo è $-25$, ovvero l'ordinata del nuovo vertice $V_1 = (5, -25)$. Cambiando di segno si ottiene allora il massimo di partenza: $\max f(x) = -(-25) = 25$.

\begin{figure}[htbp]
\centering
\begin{subfigure}[t]{0.48\textwidth}
\centering
\begin{tikzpicture}[xscale=0.45, yscale=0.085, >=Stealth]
  \useasboundingbox (-2,-33) rectangle (14,34);
  % parabola speculare (traccia)
  \draw[RuleGray, dashed, domain=-0.6:10.6, samples=60, smooth] plot (\x, {\x*\x-10*\x});
  % assi
  \draw[->] (-1.5,0) -- (11.8,0) node[right] {$x$};
  \draw[->] (0,-30) -- (0,31) node[above] {$y$};
  % curva
  \draw[Accent, thick, domain=-0.6:10.6, samples=60, smooth] plot (\x, {-\x*\x+10*\x});
  \node[Accent, font=\small, anchor=west] at (9.6,-9) {$f(x)$};
  % massimo
  \draw[dashed, TextGray] (0,25) -- (5,25) -- (5,0);
  \node[font=\scriptsize, left] at (0,25) {$25$};
  \node[font=\scriptsize, below] at (5,0) {$5$};
  \node[circle, fill=Accent, inner sep=1.8pt] at (5,25) {};
  \node[Accent, font=\small, anchor=south west] at (5.3,25.5) {$V_0$: massimo};
\end{tikzpicture}
\end{subfigure}
\hfill
\begin{subfigure}[t]{0.48\textwidth}
\centering
\begin{tikzpicture}[xscale=0.45, yscale=0.085, >=Stealth]
  \useasboundingbox (-2,-33) rectangle (14,34);
  % parabola speculare (traccia)
  \draw[RuleGray, dashed, domain=-0.6:10.6, samples=60, smooth] plot (\x, {-\x*\x+10*\x});
  % assi
  \draw[->] (-1.5,0) -- (11.8,0) node[right] {$x$};
  \draw[->] (0,-30) -- (0,31) node[above] {$y$};
  % curva
  \draw[red!70!black, thick, domain=-0.6:10.6, samples=60, smooth] plot (\x, {\x*\x-10*\x});
  \node[red!70!black, font=\small, anchor=west] at (9.6,9) {$-f(x)$};
  % minimo
  \draw[dashed, TextGray] (0,-25) -- (5,-25) -- (5,0);
  \node[font=\scriptsize, left] at (0,-25) {$-25$};
  \node[font=\scriptsize, above] at (5,0) {$5$};
  \node[circle, fill=red!70!black, inner sep=1.8pt] at (5,-25) {};
  \node[red!70!black, font=\small, anchor=north west] at (5.3,-25.5) {$V_1$: minimo};
\end{tikzpicture}
\end{subfigure}
\caption{Una parabola e la sua opposta: a sinistra $f(x) = -x^2+10x$ (in verde petrolio) con il massimo nel vertice $V_0 = (5,25)$, a destra $-f(x) = x^2-10x$ (in rosso) con il minimo nel vertice $V_1 = (5,-25)$. In entrambi i grafici la curva tratteggiata grigia è l'altra parabola: cambiare di segno la funzione equivale a ribaltarla rispetto all'asse $x$, e trasforma il massimo in un minimo di valore opposto.}
\label{fig:parabole}
\end{figure}

\section{Classificazione dei problemi}
I problemi di ottimizzazione si distinguono in due grandi classi, a seconda di come è fatta la regione ammissibile $X$.

\begin{definizione}[Problema non vincolato]
Un problema di ottimizzazione è \textbf{non vincolato} se la regione ammissibile coincide con tutto lo spazio, ovvero $X = \mathbb{R}^n$:
\[
\min_{x \in \mathbb{R}^n} f(x).
\]
\end{definizione}

In questo caso non ci sono vincoli sulle variabili. Questi problemi sono sempre e solo non lineari: una funzione lineare non nulla, infatti, è illimitata inferiormente su tutto $\mathbb{R}^n$ e quindi non ammette un punto di minimo.

\paragraph{Esempio.}
Cerchiamo il minimo di
\[
\min_{x \in \mathbb{R}} \, x^2 - 10x.
\]
\textbf{Soluzione.} Si tratta della funzione $-f(x)$ dell'esempio precedente, ovvero di una parabola con vertice $V_1 = (5, -25)$. Non essendoci vincoli, l'ottimo coincide con il vertice: $x^* = 5$, con valore della funzione obiettivo pari a $-25$.

\begin{definizione}[Problema vincolato]
Un problema di ottimizzazione è \textbf{vincolato} se la regione ammissibile è descritta da vincoli di disuguaglianza e/o di uguaglianza sulle variabili:
\[
X = \{x \in \mathbb{R}^n : g_i(x) \le 0,\ i = 1, \dots, m;\ h_j(x) = 0,\ j = 1, \dots, p\},
\]
ovvero il problema si scrive come
\[
\begin{aligned}
&\min f(x) \\
&g_i(x) \le 0, \quad i = 1, \dots, m \\
&h_j(x) = 0, \quad j = 1, \dots, p
\end{aligned}
\]
dove $g_i : \mathbb{R}^n \to \mathbb{R}$ sono i $m$ vincoli di disuguaglianza e $h_j : \mathbb{R}^n \to \mathbb{R}$ sono i $p$ vincoli di uguaglianza. Un punto è quindi ammissibile solo se li rispetta tutti\footnote{Da notare che la $i$ e la $j$ indicano rispettivamente il numero di vincoli di disuguaglianza e di uguaglianza.} contemporaneamente.
\end{definizione}

\paragraph{Esempio.}
Aggiungiamo al problema precedente un vincolo:
\[
\begin{aligned}
&\min \, x^2 - 10x \\
&x - 3 \le 0
\end{aligned}
\]
ovvero $g_1(x) = x - 3$, con $m = 1$ e $p = 0$ (c'è un solo vincolo, di disuguaglianza). Il vincolo si legge come una restrizione sulla scelta di $x$: $x - 3 \le 0$ equivale a $x \le 3$, per cui la regione ammissibile è $X = (-\infty, 3]$. La funzione obiettivo è la stessa di prima, ma ora si può scegliere $x$ solo in $X$, come mostrato nella figura~\ref{fig:vincolato}.

\begin{figure}[htbp]
\centering
\begin{tikzpicture}[xscale=0.6, yscale=0.11, >=Stealth]
  \useasboundingbox (-3.6,-36) rectangle (12.6,17);
  % regione ammissibile x <= 3
  \fill[Accent!10] (-3,-34) rectangle (3,14);
  \node[Accent, font=\small, anchor=south] at (0,-34) {$X=(-\infty,3]$};
  % assi
  \draw[->] (-3,0) -- (12,0) node[right] {$x$};
  \draw[->] (0,-13) -- (0,15) node[above] {$y$};
  \draw[Accent, very thick] (-3,0) -- (3,0);
  % bordo del vincolo
  \draw[Accent, thick] (3,-34) -- (3,14);
  \node[Accent, font=\small, anchor=west] at (3.2,10) {$x = 3$};
  % parabola: tratto ammissibile e tratto escluso dal vincolo
  \draw[red!70!black, thick, domain=-1:3, samples=50, smooth] plot (\x, {\x*\x-10*\x});
  \draw[TextGray, dashed, domain=3:11, samples=50, smooth] plot (\x, {\x*\x-10*\x});
  \node[red!70!black, font=\small, anchor=east] at (-1.1,11) {$f(x)$};
  % tacche
  \draw (3,0.8) -- (3,-0.8);
  \draw (5,0.8) -- (5,-0.8);
  \node[font=\scriptsize, above left] at (3,0.8) {$3$};
  \node[font=\scriptsize, above] at (5,0.8) {$5$};
  % minimo non vincolato (non ammissibile)
  \node[circle, draw=TextGray, fill=white, inner sep=1.6pt] at (5,-25) {};
  \node[TextGray, font=\small, anchor=north] at (5,-26.5) {$V=(5,-25)$};
  % minimo vincolato
  \node[circle, fill=red!70!black, inner sep=1.8pt] at (3,-21) {};
  \node[red!70!black, font=\small, anchor=east] at (2.7,-24) {$x^*=3,\ f(x^*)=-21$};
\end{tikzpicture}
\caption{Problema $\min x^2-10x$ con vincolo $x \le 3$: la zona colorata è la regione ammissibile $X = (-\infty, 3]$, la curva rossa è la funzione obiettivo sui punti ammissibili e il tratto grigio tratteggiato è la parte esclusa dal vincolo. Il vertice $V = (5,-25)$, minimo del problema non vincolato, non è ammissibile; l'ottimo del problema vincolato è il punto pieno $x^*=3$, sul bordo di $X$, con valore $-21$.}
\label{fig:vincolato}
\end{figure}

\subparagraph{Soluzione.} Procediamo in tre passi:
\begin{enumerate}
    \item Senza il vincolo il minimo sarebbe nel vertice $x = 5$, con valore $-25$.
    \item Questo punto però non è ammissibile, perché $5 > 3$ e quindi non rispetta il vincolo.
    \item Poiché la parabola è decrescente per $x < 5$, all'interno di $X$ la funzione obiettivo diminuisce man mano che $x$ si avvicina a $3$: il valore più basso si ottiene quindi nel punto ammissibile più vicino al vertice, cioè $x^* = 3$, con $f(3) = 9 - 30 = -21$.
\end{enumerate}
Il vincolo ha quindi spostato l'ottimo e ne ha peggiorato il valore, che passa da $-25$ a $-21$. Si noti che $x^*$ si trova sul bordo della regione ammissibile, dove il vincolo è verificato con uguaglianza.

Tra i problemi vincolati ritroviamo la \textbf{programmazione lineare} (PL), cioè il caso in cui $f$ e tutte le funzioni dei vincoli sono lineari, la \textbf{programmazione lineare intera} (PLI), in cui le variabili devono inoltre essere intere, e la \textbf{programmazione non lineare}, in cui almeno una tra $f$, $g_i$ e $h_j$ non è lineare.

\section{Definizioni}\label{sec:definizioni}
Dato il problema $\min_{x \in X} f(x)$, valgono le seguenti definizioni.

\begin{definizione}[Problema inammissibile]
Un problema di ottimizzazione è \textbf{inammissibile} se $X = \emptyset$, ovvero se nessun punto rispetta tutti i vincoli.
\end{definizione}

\begin{esempiofig}
Consideriamo ancora la funzione $x^2 - 10x$, ma con due vincoli: $x \le 3$ e $x \ge 5$. Un punto ammissibile dovrebbe rispettarli entrambi, ma nessun numero è contemporaneamente minore o uguale a $3$ e maggiore o uguale a $5$: come mostra il grafico a destra, le due zone colorate non hanno punti in comune. Quindi $X = \emptyset$ e il problema è inammissibile. In questo caso non ha senso cercare un minimo, perché non esiste alcuna scelta possibile.
\tcblower
\resizebox{\linewidth}{!}{%
\begin{tikzpicture}[xscale=0.5, yscale=0.09, >=Stealth]
  \useasboundingbox (-1.8,-30) rectangle (12.6,16);
  \fill[Accent!14] (-1,-28) rectangle (3,12);
  \fill[red!14] (5,-28) rectangle (11,12);
  \draw[->] (-1,0) -- (11.8,0) node[right, font=\scriptsize] {$x$};
  \draw[->] (0,-28) -- (0,14) node[above, font=\scriptsize] {$y$};
  \draw[TextGray, dashed, domain=-1:11, samples=50, smooth] plot (\x, {\x*\x-10*\x});
  \draw[Accent, thick] (3,-28) -- (3,12);
  \draw[red!70!black, thick] (5,-28) -- (5,12);
  \node[Accent, font=\scriptsize] at (1,8) {$x\le 3$};
  \node[red!70!black, font=\scriptsize] at (8,8) {$x\ge 5$};
  \node[font=\scriptsize, fill=white, inner sep=1.5pt] at (4,-10) {$X=\emptyset$};
\end{tikzpicture}
}
\end{esempiofig}

\begin{definizione}[Problema illimitato]
Un problema di ottimizzazione è \textbf{illimitato} se, comunque si scelga $k > 0$, esiste un punto $x \in X$ tale che $f(x) < -k$, ovvero se $f$ non è limitata inferiormente su $X$.
\end{definizione}

\begin{esempiofig}
Consideriamo il problema
\[
\min_{x \in \mathbb{R}} \, -x^2 + 10x.
\]
Per qualsiasi $k > 0$ vale $f(x) < -k$ quando $x^2 - 10x - k > 0$, e questa disuguaglianza è verificata da tutti gli $x$ abbastanza grandi in valore assoluto. Ad esempio, per $k = 20$ si ha $f(13) = -169 + 130 = -39 < -20$. La funzione obiettivo scende quindi sotto qualsiasi valore fissato, come si vede dal grafico a destra: il problema è illimitato e non ammette un minimo.
\tcblower
\resizebox{\linewidth}{!}{%
\begin{tikzpicture}[xscale=0.36, yscale=0.06, >=Stealth]
  \useasboundingbox (-4.2,-52) rectangle (15.5,32);
  \draw[->] (-4,0) -- (14.8,0) node[right, font=\scriptsize] {$x$};
  \draw[->] (0,-50) -- (0,30) node[above, font=\scriptsize] {$y$};
  \draw[Accent, thick, domain=-3.4:13.4, samples=60, smooth] plot (\x, {-\x*\x+10*\x});
  \node[Accent, font=\scriptsize, anchor=west] at (9,22) {$f(x)$};
  \draw[dashed, TextGray] (-4,-20) -- (14.5,-20);
  \node[TextGray, font=\scriptsize, anchor=south west] at (0.2,-19) {$y=-k$};
  \node[circle, fill=red!70!black, inner sep=1.5pt] at (13,-39) {};
  \node[red!70!black, font=\scriptsize, anchor=east] at (12.3,-42) {$f(x)<-k$};
\end{tikzpicture}
}
\end{esempiofig}

\begin{definizione}[Ottimo globale]
Un punto $x^* \in X$ è un \textbf{ottimo globale} se $f(x^*) \le f(x)$ per ogni $x \in X$.
\end{definizione}

\begin{esempiofig}
Riprendiamo il problema non vincolato $\min_{x \in \mathbb{R}} \, x^2 - 10x$. Il punto $x^* = 5$ è un ottimo globale, perché per ogni $x \in \mathbb{R}$ si ha
\[
f(x) - f(5) = x^2 - 10x + 25 = (x - 5)^2 \ge 0
\]
ovvero $f(x) \ge -25 = f(5)$: nessun punto può avere un valore della funzione obiettivo inferiore a quello di $x^*$, come si vede nel grafico a destra dove la curva non scende mai sotto la retta tratteggiata.
\tcblower
\resizebox{\linewidth}{!}{%
\begin{tikzpicture}[xscale=0.5, yscale=0.09, >=Stealth]
  \useasboundingbox (-1.8,-33) rectangle (12.6,16);
  \draw[->] (-1,0) -- (11.8,0) node[right, font=\scriptsize] {$x$};
  \draw[->] (0,-28) -- (0,14) node[above, font=\scriptsize] {$y$};
  \draw[Accent, thick, domain=-0.9:10.9, samples=50, smooth] plot (\x, {\x*\x-10*\x});
  \node[Accent, font=\scriptsize, anchor=west] at (9.2,9) {$f(x)$};
  \draw[dashed, TextGray] (-1,-25) -- (11,-25);
  \draw[dashed, TextGray] (5,0) -- (5,-25);
  \node[font=\scriptsize, above] at (5,0.5) {$5$};
  \node[circle, fill=red!70!black, inner sep=1.6pt] at (5,-25) {};
  \node[red!70!black, font=\scriptsize, anchor=north] at (5,-27) {$x^*=5$};
  \node[TextGray, font=\scriptsize, anchor=north west] at (6.2,-26.5) {$f(x^*)=-25$};
\end{tikzpicture}
}
\end{esempiofig}

\begin{definizione}[Ottimo locale]
Un punto $x^* \in X$ è un \textbf{ottimo locale} se esiste $\varepsilon > 0$ tale che $f(x^*) \le f(x)$ per ogni $x \in X$ con $\|x - x^*\| < \varepsilon$, ovvero se è un minimo all'interno di un intorno di $x^*$. L'ottimo locale è \textbf{stretto} se $f(x^*) < f(x)$ per ogni $x \neq x^*$ dell'intorno, cioè se $x^*$ è l'unico punto di minimo in quell'intorno.
\end{definizione}

\begin{esempiofig}
Consideriamo il problema $\min f(x) = x^3 - 3x$ con $X = [-3, 3]$. Nel punto $x = 1$ vale
\[
f(x) - f(1) = x^3 - 3x + 2 = (x - 1)^2 (x + 2)
\]
e per ogni $x$ nell'intorno $|x - 1| < 1$, ovvero in $(0, 2)$, il fattore $x + 2$ è positivo. Di conseguenza $f(x) > f(1) = -2$ per ogni $x \neq 1$ dell'intorno: $x^* = 1$ è un ottimo locale stretto. Non è però un ottimo globale, perché $f(-3) = -18 < -2$. L'ottimo globale del problema è infatti $x = -3$, estremo della regione ammissibile. Come si vede dal grafico a destra, il minimo locale è il punto più basso solo nei dintorni di $x = 1$, mentre altrove su $X$ la funzione scende ancora di più.
\tcblower
\resizebox{\linewidth}{!}{%
\begin{tikzpicture}[xscale=0.85, yscale=0.14, >=Stealth]
  \useasboundingbox (-3.6,-22) rectangle (4.3,22);
  \draw[->] (-3.3,0) -- (3.9,0) node[right, font=\scriptsize] {$x$};
  \draw[->] (0,-21) -- (0,21) node[above, font=\scriptsize] {$y$};
  \draw[Accent, thick, domain=-3:3, samples=60, smooth] plot (\x, {\x*\x*\x-3*\x});
  \draw[red!70!black, very thick, domain=0:2, samples=30, smooth] plot (\x, {\x*\x*\x-3*\x});
  \draw[red!70!black, very thick] (0,0) -- (2,0);
  \node[red!70!black, font=\scriptsize, anchor=south] at (1,0.8) {$\varepsilon=1$};
  \node[circle, fill=red!70!black, inner sep=1.6pt] at (1,-2) {};
  \node[red!70!black, font=\scriptsize, anchor=north] at (1.2,-4) {min locale};
  \node[circle, fill=Accent, inner sep=1.6pt] at (-3,-18) {};
  \node[Accent, font=\scriptsize, anchor=west] at (-2.75,-19) {min globale};
\end{tikzpicture}
}
\end{esempiofig}

Un ottimo locale non stretto è invece un punto che appartiene a una regione confinata di punti con lo stesso valore minimo di $f$. Naturalmente ogni ottimo globale è anche un ottimo locale, ma non è detto il contrario.

\begin{osservazione}
Nella programmazione lineare la distinzione tra ottimo globale e ottimo locale non esiste: ogni ottimo locale è anche un ottimo globale (si veda la sezione~\ref{sec:graf}).
\end{osservazione}

\section{Esistenza della soluzione}
Assegnati $X$ ed $f$, potrebbero anche non esistere soluzioni ottime. Quando si cerca $\min_{x \in X} f(x)$ possono verificarsi quattro casi:
\begin{enumerate}
    \item La regione ammissibile è vuota: il problema è inammissibile e non c'è alcuna scelta possibile.
    \item La regione ammissibile non è vuota, ma $f$ non è limitata inferiormente: il problema è illimitato e non esiste un minimo.
    \item La regione ammissibile non è vuota e $f$ è limitata inferiormente, ma non esistono punti di minimo globale.
    \item La regione ammissibile non è vuota ed esiste un punto di minimo globale; in questo caso $f$ è ovviamente limitata inferiormente.
\end{enumerate}

Risolvere un problema di ottimizzazione significa quindi, in pratica:
\begin{itemize}
    \item Stabilire se l'insieme ammissibile è non vuoto, oppure concludere che non esistono soluzioni ammissibili.
    \item Stabilire se esistono soluzioni ottime, oppure dimostrare che il problema non ammette soluzioni ottime.
    \item Determinare, eventualmente in modo approssimato, una soluzione ottima.
\end{itemize}

\section*{Riferimenti}
Il materiale visto in questo capitolo comprende:
\begin{itemize}
    \item Materiale visto a lezione (2 ottobre 2026).
    \item Dispensa del corso \emph{Introduzione all'ottimizzazione}, capitoli 1--3.
\end{itemize}
